Selecting locations and ordering their visits are coupled decisions in the
prize-collecting covering tour problem (PCCTP). A tour must include mandatory
locations, serve demands through nearby visited locations and collect a
minimum prize. An optional location can contribute to both coverage and
prize collection, while also changing the routing cost. Consequently, a
shorter ordering of a fixed visited set and a better choice of that set are
different sources of improvement. An algorithm that changes visits must
account for the feasibility requirements that those visits were satisfying.

Existing PCCTP methods already address this interaction through several
optimization components. Reactive GRASP with GENIUS construction and variable
neighborhood descent provides a heuristic precedent~\cite{climaco2021reactive}.
Reduction rules, valid inequalities, branch-and-cut and a heuristic--integer
programming hybrid have also been studied~\cite{climaco2023approaches}.
The later MDM-GRASP approach uses patterns mined from high-quality
solutions~\cite{climaco2026mdm}. These precedents make a comparison with
construction alone insufficient and prevent treating generic hybridization
as a contribution by itself. The relevant question is whether a specified
combination improves the solution process after its additional work is charged.

This study evaluates destruction and repair after MDM initialization, with
route ordering performed by LKH and with greedy and restricted integer-repair
alternatives examined separately. It asks three questions. Does changing the
visited set add value beyond reordering the same set from common starts?
Does joint integer insertion improve on greedy repair after matched
destructions? Finally, does a complete hybrid improve on a sufficiently
developed MDM reference under an explicit total-time comparison? Separating
these questions avoids attributing a strong ordering component's gains to
repair, or presenting a conditional improvement from an existing solution
as a complete algorithmic advantage.

The completed evidence is exploratory. It includes a positive short-budget
training comparison, a failed validation screen, a conditional fixed-support
replay and an unfavorable mean result for the evaluated joint-insertion MIP
variant. The manuscript reports these outcomes with their denominators,
geometry-level sensitivity and independent feasibility checks. The
longer-horizon comparison and the iteration-limited reference study remain
necessary before making a broader performance claim. The formulation,
neighborhood descriptions and audited exploratory results presented here
therefore do not establish a new best-known solution, a novel general
hybridization principle or superiority over the published PCCTP methods.
